\section{Conclusion}
\label{sec:conclusion}

This paper estimates the fiscal costs of judicial enforcement in public health procurement, exploiting granular bid-level data on pharmaceutical purchases in S\~ao Paulo, Brazil. Our findings document that court-mandated urgent procurement significantly distorts procurement outcomes: reference prices rise by 2.7\%, negotiated prices increase by 3.1--5.4\%, and bidder participation falls by 4.1--5.5\%, all within item, year, and buyer comparisons.

The ``under the gun'' analysis reveals that the sanction channel is particularly consequential. Comparing litigated and administrative urgent purchases---which share identical planning constraints but differ in penalty exposure---we find that judicial sanctions alone account for a 23--30\% price premium. This magnitude suggests that the threat of judicial punishment fundamentally alters the bargaining environment, distorting procurement incentives well beyond what planning disruption alone would imply.

The finding that urgent tenders are more likely to succeed, despite worse terms, is consistent with the mechanism: procurement officials under judicial pressure prioritize compliance over cost-efficiency, accepting higher prices to avoid the severe consequences of tender failure. This pattern echoes the broader insight from the accountability literature that oversight mechanisms can generate unintended costs when they distort rather than align incentives \citep{rasul2018management, duflo2018obligations}.

Our results carry several policy implications. First, the large cost premium attributable to judicial sanctions suggests that revising penalty mechanisms could yield significant fiscal savings. More proportionate sanctions---or institutional arrangements that shield procurement officials from personal liability for good-faith efforts---could preserve accountability while reducing the distortionary ``under the gun'' effect. Second, enhanced coordination between the judiciary and the executive, such as expanding the administrative request mechanism, could reduce the volume of litigated purchases and the associated costs. Third, the concentration of enforcement costs in markets with less competition suggests that targeted supply-side interventions---such as framework agreements or strategic stockpiling for commonly litigated items---could mitigate the price premium.

More broadly, our findings highlight a tension at the heart of right-to-health enforcement: judicial mandates intended to guarantee individual access to health care impose substantial costs on the public budget, potentially reducing the resources available for the broader health system. The estimated price premiums, applied to the approximately US\$300~million spent annually on litigated health items in S\~ao Paulo alone, imply tens of millions of dollars in excess costs. The geographic evidence presented in Section~\ref{sec:institutional} and in the Appendix confirms that these costs are concentrated in specific municipalities and PBUs, suggesting that targeted interventions could be particularly effective. Institutional designs that balance individual rights with procurement efficiency---rather than treating them as competing objectives---could improve welfare for both litigants and the broader population.

Several limitations suggest directions for future research. Our sample covers a single state's procurement system, and the generalizability of the magnitudes to other jurisdictions and institutional contexts remains an open question. The ``under the gun'' identification relies on the assumption that administrative and litigated purchases differ only in the penalty regime; while institutional evidence strongly supports this assumption, unobserved differences in the composition of requests cannot be entirely ruled out. Future work could exploit administrative reforms or policy changes that alter the penalty regime to provide complementary identification. Cross-country comparisons of judicial enforcement in procurement---particularly in other Latin American countries with active health litigation---would also be valuable.
